
\chapter{Vladimir Drinfeld (2018)}
\index{Drinfeld, Vladimir}

\begin{quote}
\it ``Vladimir Drinfeld's work has opened entirely new horizons in mathematics. From Drinfeld modules and the Langlands correspondence for function fields to quantum groups and the geometric Langlands program, his ideas have created new fields and transformed existing ones.'' --- Wolf Prize Citation, 2018
\end{quote}

\section{Main Contributions}

Vladimir Drinfeld (born 1954) is a Ukrainian--American mathematician whose work has revolutionized number theory, representation theory, and mathematical physics. He was awarded the Fields Medal in 1990 and the Wolf Prize in Mathematics in 2018 alongside Alexander Beilinson.

The Wolf Prize citation reads: ``for their contributions to algebraic geometry, representation theory, and number theory.''

Drinfeld's core contributions include:
\begin{itemize}
    \item \textbf{Drinfeld Modules:} The introduction of elliptic modules (now called Drinfeld modules), which are function field analogues of elliptic curves. These are central to the arithmetic of function fields.
    \item \textbf{Langlands Correspondence for Function Fields:} A proof of the Langlands correspondence for $\GL(2)$ over global function fields, a landmark result connecting Galois representations to automorphic forms.
    \item \textbf{Quantum Groups:} The discovery (independently by Jimbo) of quantum groups, which are deformations of universal enveloping algebras of Lie algebras. These are the mathematical foundation of the quantum inverse scattering method.
    \item \textbf{Geometric Langlands Program:} With Beilinson, the development of the geometric Langlands program, including the quantization of Hitchin's integrable system.
    \item \textbf{Drinfeld--Sokolov Reduction:} A powerful technique relating Kac--Moody algebras to integrable systems.
    \item \textbf{The Quantum Yang--Baxter Equation:} The classification of solutions to the quantum Yang--Baxter equation via Drinfeld's theory of Hopf algebras.
\end{itemize}

\section{Origin of the Problem}

\subsection{The Langlands Program over Function Fields}

The classical Langlands program concerns number fields. Drinfeld realized that the function field case (fields of the form $\F_q(C)$ for a curve $C$ over a finite field) might be more tractable while still capturing the essential structure.

The key advantage: in function fields, there is an additional geometric structure (the curve $C$) that can be exploited using algebraic geometry.

\subsection{Drinfeld Modules}

Elliptic curves over number fields are central to modern number theory. Over function fields, the analogue is not an elliptic curve but a new object that Drinfeld introduced: a module over the ring $\F_q[t]$ whose endomorphism ring is a field of characteristic $p$.

\begin{definition}[Drinfeld Module]
A \textbf{Drinfeld module} of rank $r$ over a field $K$ of characteristic $p$ is an $\F_q[t]$-module structure on the additive group $\mathbb{G}_a(K)$ such that the action of $t$ is given by a polynomial:
\[
t \cdot x = a_0 x^q + a_1 x^{q^2} + \cdots + a_r x^{q^r}
\]
with $a_0 \neq 0$.
\end{definition}

\subsection{Quantum Groups}

The quantum inverse scattering method in statistical mechanics led to the Yang--Baxter equation. The question: what algebraic structure underlies solutions to this equation? Drinfeld discovered that the answer is a deformation of the universal enveloping algebra $U(\mathfrak{g})$, which he called a quantum group.

\section{How It Was Solved and the Intuitions/Tools Needed}

\subsection{The Rank 2 Langlands Correspondence}

Drinfeld proved the Langlands correspondence for $\GL(2)$ over function fields by:
\begin{enumerate}
    \item Constructing the Drinfeld modular variety (a moduli space of Drinfeld modules).
    \item Studying the action of the Hecke algebra on the cohomology of this variety.
    \item Using the Eichler--Shimura congruence relations to link the Hecke action to the Frobenius action on Galois representations.
\end{enumerate}

This approach was conceptually parallel to the proof of the Shimura--Taniyama--Weil conjecture (modularity of elliptic curves) but in the function field context.

\subsection{Quantum Groups}

Drinfeld's construction of quantum groups starts with the Cartan data of a semisimple Lie algebra $\mathfrak{g}$ and produces a Hopf algebra $U_q(\mathfrak{g})$ over $\C(q)$ satisfying:
\begin{itemize}
    \item $U_q(\mathfrak{g})$ is a deformation of $U(\mathfrak{g})$: as $q \to 1$, $U_q(\mathfrak{g}) \to U(\mathfrak{g})$.
    \item $U_q(\mathfrak{g})$ has a universal R-matrix satisfying the Yang--Baxter equation.
    \item The representation theory of $U_q(\mathfrak{g})$ is a $q$-deformation of the classical representation theory.
\end{itemize}

\section{Mathematical Theory}

\subsection{Drinfeld Modules}

\begin{definition}[Function Field]
Let $\F_q$ be a finite field. A \textbf{function field} is a finite extension of $\F_q(t)$. It is the field of rational functions on a smooth projective curve $C$ over $\F_q$.
\end{definition}

\begin{definition}[Drinfeld Module]
Let $A = \F_q[t]$ and let $K$ be a field of characteristic $p$ with an $\F_q$-algebra homomorphism $\iota: A \to K$. A \textbf{Drinfeld module} of rank $r$ over $K$ is an $\F_q$-algebra homomorphism:
\[
\phi: A \to K\{\tau\}, \quad \phi_t = a_0 \tau^0 + a_1 \tau + \cdots + a_r \tau^r
\]
where $K\{\tau\}$ is the noncommutative ring of polynomials in the Frobenius $\tau(x) = x^q$, with the commutation $\tau a = a^q \tau$ for $a \in K$.
\end{definition}

\subsection{Langlands Correspondence for $\GL(2)$ over Function Fields}

\begin{theorem}[Drinfeld, 1974]
Let $F$ be a global function field of characteristic $p$ and let $\GL(2, \A_F)$ be the adele group. There is a natural bijection between:
\begin{enumerate}
    \item Irreducible cuspidal automorphic representations of $\GL(2, \A_F)$.
    \item Irreducible $2$-dimensional representations of the absolute Galois group $\Gal(F^{\mathrm{sep}}/F)$.
\end{enumerate}
\end{theorem}

This was the first complete proof of the Langlands correspondence for a group beyond $\GL(1)$ (class field theory).

\subsection{Quantum Groups}

\begin{definition}[Quantum Group]
Let $\mathfrak{g}$ be a semisimple Lie algebra with Cartan matrix $(a_{ij})$. The \textbf{quantum group} $U_q(\mathfrak{g})$ is the Hopf algebra over $\C(q)$ generated by $E_i, F_i, K_i^{\pm 1}$ with relations:
\[
K_i K_j = K_j K_i, \quad K_i E_j K_i^{-1} = q^{a_{ij}} E_j, \quad K_i F_j K_i^{-1} = q^{-a_{ij}} F_j
\]
\[
E_i F_j - F_j E_i = \delta_{ij} \frac{K_i - K_i^{-1}}{q - q^{-1}}
\]
and the quantum Serre relations:
\[
\sum_{k=0}^{1-a_{ij}} (-1)^k \left[\begin{matrix} 1-a_{ij} \\ k \end{matrix}\right]_q E_i^{1-a_{ij}-k} E_j E_i^k = 0
\]
\end{definition}

\begin{definition}[Universal R-Matrix]
A \textbf{universal R-matrix} for a Hopf algebra $H$ is an invertible element $R \in H \otimes H$ such that:
\[
\Delta^{\mathrm{op}}(x) = R \Delta(x) R^{-1} \quad \text{for all } x \in H
\]
and $R$ satisfies the Yang--Baxter equation:
\[
R_{12} R_{13} R_{23} = R_{23} R_{13} R_{12}
\]
\end{definition}

Every quantum group $U_q(\mathfrak{g})$ possesses a universal R-matrix, providing a systematic source of solutions to the Yang--Baxter equation.

\subsection{Drinfeld--Sokolov Reduction}

\begin{definition}[Drinfeld--Sokolov Reduction]
A procedure that associates to an affine Kac--Moody algebra a classical integrable system. Starting from a principal $\mathfrak{sl}_2$ embedding into $\mathfrak{g}$, one imposes constraints on the affine Kac--Moody algebra to obtain the Gelfand--Dickey hierarchy (including KdV, modified KdV, etc.).
\end{definition}

\section{Impact on Science and Technology}

\subsection{In Mathematics}

\begin{enumerate}
    \item \textbf{Number Theory:} Drinfeld modules are central to the arithmetic of function fields, playing the role of elliptic curves.
    \item \textbf{Representation Theory:} Quantum groups are fundamental objects, connecting Lie theory to low-dimensional topology and knot theory.
    \item \textbf{Geometric Langlands:} Drinfeld's work, with Beilinson, opened the geometric Langlands program.
    \item \textbf{Low-Dimensional Topology:} Quantum groups provide invariants of knots and 3-manifolds (Jones polynomial, Reshetikhin--Turaev invariants).
\end{enumerate}

\subsection{In Physics}

\begin{enumerate}
    \item \textbf{Statistical Mechanics:} Quantum groups solve integrable lattice models via the quantum Yang--Baxter equation.
    \item \textbf{Conformal Field Theory:} Drinfeld--Sokolov reduction relates Kac--Moody algebras to W-algebras and conformal field theories.
    \item \textbf{Quantum Field Theory:} Universal R-matrices appear in the study of quantum integrable systems.
\end{enumerate}

\section{Five Frequently Asked Questions}

\subsection*{Q1: What is a Drinfeld module?}
A Drinfeld module is a function field analogue of an elliptic curve. It is an algebraic object that allows the ring $\F_q[t]$ to act on the additive group in a nontrivial way, providing the foundation for the arithmetic of function fields.

\subsection*{Q2: What are quantum groups?}
Quantum groups are deformations of universal enveloping algebras of Lie algebras. Despite the name, they are not groups but Hopf algebras. They provide solutions to the Yang--Baxter equation and are central to mathematical physics.

\subsection*{Q3: What did Drinfeld prove about Langlands?}
He proved the Langlands correspondence for $\GL(2)$ over global function fields, establishing a bijection between automorphic representations and 2-dimensional Galois representations. This was the first proof for a non-abelian group.

\subsection*{Q4: What is the quantum Yang--Baxter equation?}
An equation $R_{12} R_{13} R_{23} = R_{23} R_{13} R_{12}$ for an operator $R$ on $V \otimes V$. It is the master equation for integrable quantum systems and the foundation of the theory of quantum groups.

\subsection*{Q5: What should an undergraduate study to understand Drinfeld's work?}
Algebraic geometry (curves, schemes), number theory (local fields, Galois theory), and representation theory (Lie algebras, Hopf algebras). Recommended: \textit{Algebraic Number Theory} (Neukirch), \textit{Introduction to Quantum Groups} (Kassel), and \textit{Drinfeld Modules} (Goss).

\section*{Exercises for the Reader}
\addcontentsline{toc}{section}{Exercises}

\begin{enumerate}
    \item [B] Show that a Drinfeld module of rank 1 corresponds to the Carlitz module.
    \item [I] Verify that the quantum group $U_q(\mathfrak{sl}_2)$ satisfies the Hopf algebra axioms.
    \item [A] Compute the universal R-matrix for $U_q(\mathfrak{sl}_2)$ explicitly.
    \item [A] Show that the Drinfeld modular variety parametrizes rank 2 Drinfeld modules.
    \item [A] Derive the KdV hierarchy from the Drinfeld--Sokolov reduction of $\widehat{\mathfrak{sl}}_2$.
    \item [I] Prove that the Langlands correspondence for $\GL(1)$ over function fields is class field theory.
    \item [I] Compute the invariant of the trefoil knot using the quantum group $U_q(\mathfrak{sl}_2)$.
\end{enumerate}

\section*{Timeline}
\addcontentsline{toc}{section}{Timeline}

\begin{tabularx}{\linewidth}{|l|X|}
\hline
\textbf{Year} & \textbf{Event} \\ \hline
1954 & Born in Kharkiv, Ukrainian SSR \\ \hline
1969 & Wins gold at the International Mathematical Olympiad \\ \hline
1974 & Proves Langlands correspondence for $\GL(2)$ over function fields \\ \hline
1978 | Ph.D.\ from Moscow State University (advisor: Yuri Manin) \\ \hline
1980 | Discovers Drinfeld modules \\ \hline
1985 | Discovers quantum groups \\ \hline
1988 | Moves to the Institute for Low Temperature Physics, Kharkiv \\ \hline
1990 | Fields Medal at ICM Kyoto \\ \hline
1998 | Moves to the University of Chicago \\ \hline
2000s | Geometric Langlands program (with Beilinson) \\ \hline
2018 | Wolf Prize in Mathematics \\ \hline
\end{tabularx}

\section*{Glossary}
\addcontentsline{toc}{section}{Glossary}

\begin{description}
    \item[Drinfeld module] A function field analogue of an elliptic curve, giving an $\F_q[t]$-module structure on $\mathbb{G}_a$.
    \item[Function field] A finite extension of $\F_q(t)$, the field of rational functions on a curve.
    \item[Geometric Langlands] A geometric analogue of the Langlands program over function fields.
    \item[Hopf algebra] An algebra with a coproduct, counit, and antipode, generalizing group algebras.
    \item[Quantum group] A deformation $U_q(\mathfrak{g})$ of the universal enveloping algebra of a Lie algebra.
    \item[Universal R-matrix] An element of $H \otimes H$ satisfying the Yang--Baxter equation.
    \item[Yang--Baxter equation] The fundamental equation $R_{12}R_{13}R_{23} = R_{23}R_{13}R_{12}$.
\end{description}

\section{Citations}

\subsection*{Primary Sources}
\begin{enumerate}
    \item V.\ Drinfeld, \textit{Elliptic modules}, Math.\ USSR Sbornik 23 (1974), 561--592.
    \item V.\ Drinfeld, \textit{Two-dimensional l-adic representations of the Galois group of a global field of characteristic $p$ and automorphic forms on $\GL(2)$}, J.\ Soviet Math.\ 36 (1987), 105--124.
    \item V.\ Drinfeld, \textit{Quantum groups}, Proc.\ ICM Berkeley 1986, 798--820.
    \item V.\ Drinfeld, \textit{Hopf algebras and the quantum Yang--Baxter equation}, Soviet Math.\ Dokl.\ 32 (1985), 254--258.
    \item V.\ Drinfeld and V.\ Sokolov, \textit{Lie algebras and equations of Korteweg--de Vries type}, J.\ Soviet Math.\ 30 (1985), 1975--2036.
\end{enumerate}

\subsection*{Expository Works}
\begin{enumerate}
    \item D.\ Goss, \textit{Basic Structures of Function Field Arithmetic}, Springer, 1996.
    \item C.\ Kassel, \textit{Quantum Groups}, Springer, 1995.
    \item E.\ Frenkel, \textit{Lectures on the Langlands Program and Conformal Field Theory}, Les Houches, 2005.
    \item V.\ Chari and A.\ Pressley, \textit{A Guide to Quantum Groups}, Cambridge Univ.\ Press, 1994.
\end{enumerate}

